\documentclass[10pt,a4paper]{amsart}
\usepackage[margin=19mm]{geometry}
\usepackage{amsmath,amssymb,mathtools}
\usepackage[scr=rsfs,cal=euler]{mathalfa}
\usepackage{xfrac}
\usepackage[unicode,hidelinks,bookmarks=false]{hyperref}
\newcommand{\ie}{i.e.\ }
\newcommand{\cf}{cf.\ }
\newcommand{\resp}{respectively\ }
\newcommand{\Subsection}{\subsection}
\providecommand{\nobreakdash}{\nobreak\hbox{-}}
\setlength{\emergencystretch}{2em}
\multlinegap0pt
\usepackage{mdframed}
\usepackage{mdframed}
\usepackage{mdframed}
\usepackage{mdframed}
\usepackage{amssymb}
\usepackage{tikz}
\usepackage{breqn}
\usepackage{mdframed}
\usepackage{float}
\newtheorem{lemma}{Lemma}
\newtheorem{definition}{Definition}
\newtheorem{restatable}{Restatable}
\newcommand{\N}{\mathbb{N}}
\newcommand{\abag}{\mathfrak{b}}
\newcommand{\abstractdecomposition}{\tau}
\newcommand{\action}{\rho}
\newcommand{\allabstractdecompositionstreewidth}[1]{\mathsf{ITD}_{#1}}
\newcommand{\allsubterms}{\mathrm{Sub}}
\newcommand{\allwitnesses}{\mathcal{W}}
\newcommand{\awitness}{\mathsf{w}}
\newcommand{\defeq}{~\dot{=}~}
\newcommand{\domain}{dom}
\newcommand{\dpcore}{\mathsf{D}}
\newcommand{\dprefutation}{R}
\newcommand{\finalwitnessgenericcore}{\mathtt{Final}}
\newcommand{\finitepowerset}[1]{\mathcal{P}_{\mathrm{fin}}(#1)}
\newcommand{\forgetvertextype}[1]{\mathtt{ForgetVertex}\{#1\}}
\newcommand{\initialsetgeneric}{\mathtt{Leaf}}
\newcommand{\introedgegeneric}[2]{\mathtt{IntroEdge}\{#1,#2\}}
\newcommand{\introvertextype}[1]{\mathtt{IntroVertex}\{#1\}}
\newcommand{\jointype}{\mathtt{Join}}
\newcommand{\leaftype}{\mathtt{Leaf}}
\newcommand{\relabelclass}{\mathcal{F}}
\newcommand{\relabelingfunction}{f}
\newcommand{\relabelsequence}{\textsc{F}}
\newcommand{\subtermToBag}{T}
\newcommand{\topbag}[1]{\topbagname(#1)}
\newcommand{\topbagname}{B}
\newcommand{\treewidthvalue}{k}
\newcommand{\vertexone}{u}
\newcommand{\vertextwo}{v}
\newcommand{\witnessset}{S}
\title{State Canonization and Early Pruning in\\Width-Based Automated Theorem Proving}
\author{Mateus de Oliveira Oliveira, Sam Urmian}
\date{Source: arXiv:2605.11025v1 (CC BY 4.0)}
\begin{document}
\maketitle

\noindent\emph{Source and licence.} State Canonization and Early Pruning in\\Width-Based Automated Theorem Proving. Version arXiv:2605.11025v1 (CC BY 4.0), distributed by arXiv under the Creative Commons Attribution 4.0 International licence, http://creativecommons.org/licenses/by/4.0/, as recorded in the arXiv licence metadata for this exact version and shown on the abstract page https://arxiv.org/abs/2605.11025v1. Full licence notice: \texttt{licenses/arxiv-2605.11025-CC-BY-4.0.txt}.\par\smallskip

\noindent\textit{Editorial scope.} Counted unit: the article's lemma treedecexistssemiref (source0.tex lines 2001--2020). The article's complete original proof of it is delivered with it (source0.tex lines 2022--2207, one \texttt{proof} environment of 186 lines). No mathematics is written, repaired or re-derived: the statement and the proof are the article's own lines, in the article's own notation, and no retained line is substituted or re-flowed.  Retained uncounted as the source-local context the proof needs: lines 1692--1696; lines 1704--1707; lines 1777--1811; lines 1821--1831; lines 1833--1864; lines 1951--1989.  Nothing was omitted from any retained span, so the package carries no omission record.  No retained line contains a citation, so the wrapper carries no bibliography and no bibliography entry.  No retained line contains a figure environment, an image-inclusion command or drawing code, so no image is delivered and the package declares no asset dependency.  Editorial formatting only, in addition: an AMS \texttt{amsart} preamble that restates the article's own declarations and macros used by the retained lines, namely the environments \texttt{lemma}, \texttt{definition}, \texttt{restatable} on separate counters, together with the standard packages those lines need.  The retained environments print in this document as Lemma 1, Definition 1, Restatable 1 in order of appearance, whereas the article prints the same items with the numbering of its own full document.  The complete native source of the article as distributed in the arXiv e-print for this exact version is bundled unchanged as \texttt{sources/200333/source0.tex}.
\begin{lemma}[Injective extension]
\label{lemma:injective-extension}
Let $\abag\subseteq [\treewidthvalue+1]$ and let $\vertexone\in \abag$.
If $h:\abag\setminus\{\vertexone\}\to [\treewidthvalue+1]$ is injective, then there exists an injective extension $\hat h:\abag\to[\treewidthvalue+1]$ of $h$, i.e., $\hat h|_{\abag\setminus\{\vertexone\}}=h$.
\end{lemma}

\begin{lemma}[Injective image containment]
\label{lemma:injective-image-containment}
Let $f\in \relabelclass_{\treewidthvalue}$ and let $X,Y\subseteq \domain(f)$. If $f(X)\subseteq f(Y)$, then $X\subseteq Y$.
\end{lemma}

\begin{definition}[Witness Action]\label{definition:witnessaction}
Let $\dpcore$ be a DP-core and $k\in \N$. A witness action for $\dpcore[\treewidthvalue]$ is a partial function
$\action_{\dpcore}^{\treewidthvalue}$ from $\relabelclass_\treewidthvalue \times \dpcore[\treewidthvalue].\allwitnesses$ to $\dpcore[\treewidthvalue].\allwitnesses$, defined exactly on the supported pairs
\[
(\relabelingfunction,\awitness)
\quad\text{with}\quad
\mathrm{Lbl}_{\treewidthvalue}(\awitness)\subseteq \domain(\relabelingfunction).
\]
We say that $\action_{\dpcore}^{\treewidthvalue}$ is an action for $\dpcore[\treewidthvalue]$ if
the following conditions are satisfied for each $\relabelingfunction,\relabelingfunction'\in \relabelclass_{\treewidthvalue}$, and
each $\awitness \in \dpcore[\treewidthvalue].\allwitnesses$, whenever the displayed expressions are supported and hence defined.
Since $\relabelingfunction$ is injective, $\relabelingfunction^{-1}$ denotes the inverse map on its image $\relabelingfunction(\domain(\relabelingfunction))$ (so $\relabelingfunction^{-1}\in \relabelclass_\treewidthvalue$ with domain $\relabelingfunction(\domain(\relabelingfunction))$).
\begin{enumerate}
	    \item \label{condition:relabeldpcore-condition-one}
	    Finality invariance:
	    \[
	    \dpcore[\treewidthvalue].\finalwitnessgenericcore(\awitness)=1
	    \quad \text{if and only if}\quad
	    \dpcore[\treewidthvalue].\finalwitnessgenericcore(\action_{\dpcore}^{\treewidthvalue}(\relabelingfunction,\awitness))=1.
	    \]
	    \item \label{condition:relabeldpcore-condition-support-transport}
	    Label support transport:
	    \[
	    \mathrm{Lbl}_{\treewidthvalue}(\action_{\dpcore}^{\treewidthvalue}(\relabelingfunction,\awitness))
	    = \relabelingfunction(\mathrm{Lbl}_{\treewidthvalue}(\awitness)).
	    \]
	    \item Inverse relabeling: \label{condition:relabeldpcore-condition-two} \[\action_{\dpcore}^{\treewidthvalue}(\relabelingfunction^{-1},\action_{\dpcore}^{\treewidthvalue}(\relabelingfunction,\awitness)) = \awitness.\]
	    \item Composition: \label{condition:relabeldpcore-condition-three} \[\action_{\dpcore}^{\treewidthvalue}(\relabelingfunction\circ \relabelingfunction',\awitness) =
	    \action_{\dpcore}^{\treewidthvalue}(\relabelingfunction,\action_{\dpcore}^{\treewidthvalue}(\relabelingfunction',\awitness)).\]
	    \item \label{condition:relabeldpcore-condition-extension-invariance}
	    Extension invariance:
	    if $\relabelingfunction'\in \relabelclass_{\treewidthvalue}$ satisfies $\domain(\relabelingfunction)\subseteq \domain(\relabelingfunction')$ and $\relabelingfunction'|_{\domain(\relabelingfunction)}=\relabelingfunction$, then
		    \[
		    \action_{\dpcore}^{\treewidthvalue}(\relabelingfunction',\awitness)=\action_{\dpcore}^{\treewidthvalue}(\relabelingfunction,\awitness).
		    \]

			    \item DP-core Consistency Axioms:
				    \begin{enumerate}
				    \item \label{condition:relabeldpcore-condition-four}
				    For each $\vertexone\in \domain(\relabelingfunction)$,
				    \[
				    \begin{aligned}
				    \action_{\dpcore}^{\treewidthvalue}(\relabelingfunction,\introvertextype{\vertexone}(\awitness))
				    &=
				    \introvertextype{\relabelingfunction(\vertexone)}\bigl(\action_{\dpcore}^{\treewidthvalue}(\relabelingfunction,\awitness)\bigr).
				    \end{aligned}
				    \]

		    \item \label{condition:relabeldpcore-condition-five}
				    For each $\vertexone\in \domain(\relabelingfunction)$,
				    \[
				    \begin{aligned}
				    \action_{\dpcore}^{\treewidthvalue}(\relabelingfunction,\forgetvertextype{\vertexone}(\awitness))
				    &=
				    \forgetvertextype{\relabelingfunction(\vertexone)}\bigl(\action_{\dpcore}^{\treewidthvalue}(\relabelingfunction,\awitness)\bigr).
				    \end{aligned}
				    \]
		    \item \label{condition:relabeldpcore-condition-six}
				    For each two distinct $\vertexone,\vertextwo\in \domain(\relabelingfunction)$,
				    \[
				    \begin{aligned}
				    \action_{\dpcore}^{\treewidthvalue}(\relabelingfunction,\introedgegeneric{\vertexone}{\vertextwo}(\awitness))
				    &=
				    \introedgegeneric{\relabelingfunction(\vertexone)}{\relabelingfunction(\vertextwo)}\\
				    &\qquad\bigl(\action_{\dpcore}^{\treewidthvalue}(\relabelingfunction,\awitness)\bigr).
				    \end{aligned}
				    \]
	    \item \label{condition:relabeldpcore-condition-seven}
				    For each $\awitness'\in  \dpcore[\treewidthvalue].\allwitnesses$ with $\mathrm{Lbl}_{\treewidthvalue}(\awitness')\subseteq \domain(\relabelingfunction)$,
				    \[
				    \begin{aligned}
				    \action_{\dpcore}^{\treewidthvalue}(\relabelingfunction,\jointype(\awitness,\awitness'))
				    &=
				    \jointype\Bigl(\action_{\dpcore}^{\treewidthvalue}(\relabelingfunction,\awitness),\\
				    &\qquad\action_{\dpcore}^{\treewidthvalue}(\relabelingfunction,\awitness')\Bigr).
				    \end{aligned}
				    \]
			    \end{enumerate}
\end{enumerate}
\end{definition}

\begin{definition}[$\relabelsequence$-Relabeled Refutation]\label{definition:RelabeledRefutation}
Let $\relabelsequence = (\relabelingfunction_1,\dots,\relabelingfunction_m)$ be a sequence of relabeling functions in $\relabelclass_\treewidthvalue$ and $\dpcore$ be a DP-core. An $\relabelsequence$-relabeled $(\treewidthvalue,\dpcore)$-refutation is a sequence of pairs
$$\dprefutation \equiv (\abag_0,\witnessset_0)(\abag_1,\witnessset_1)\dots (\abag_m,\witnessset_m)$$ satisfying the following conditions:
\begin{enumerate}
\setlength\itemsep{0.5em}
\item \label{sequence-relabeling-condition-one}
$(\abag_0,\witnessset_0) = (\emptyset,\dpcore[\treewidthvalue].\initialsetgeneric)$.
\item \label{sequence-relabeling-condition-two}
$(\abag_m,\witnessset_m)$ is $(\treewidthvalue,\dpcore)$-inconsistent, i.e., $\witnessset_m$ has no final local witness.
	\item \label{sequence-relabeling-condition-wellformed}
	For each $i\in\{0,1,\dots,m\}$, the state $(\abag_i,\witnessset_i)$ is well-formed, i.e., $\witnessset_i\in \finitepowerset{\dpcore[\treewidthvalue].\allwitnesses}$ and $\mathrm{Lbl}_{\treewidthvalue}(\witnessset_i)\subseteq \abag_i$.
					\item \label{sequence-relabeling-condition-three}
					For each $i\in [m]$, there is some $j\in \{0\}\cup[i-1]$ such that $(\abag_i,\witnessset_i)$ is equal to one of the following pairs.
					In each alternative, $\relabelingfunction_i$ is taken with domain equal to the pre-relabeling output bag displayed in the first component: respectively $\abag_j\cup\{\vertexone\}$, $\abag_j\setminus\{\vertexone\}$, $\abag_j$, and $\abag_j$.
				\begin{enumerate}
	\item \label{sequence-relabeling-IntroVertex}
	%$(\abag_i,\witnessset_i)=$\\
	$(\relabelingfunction_i(\abag_j\cup \{\vertexone\}), \action_{\dpcore}^{\treewidthvalue}(\relabelingfunction_i,\introvertextype{\vertexone}(\witnessset_j)))$
%\\
for some $\vertexone\notin \abag_j$.
\item \label{sequence-relabeling-ForgetVertex}
%$(\abag_i,\witnessset_i) = $ \\
$(\relabelingfunction_i(\abag_j\setminus \{\vertexone\}), \action_{\dpcore}^{\treewidthvalue}(\relabelingfunction_i,\forgetvertextype{\vertexone}(\witnessset_j)))$
%\\
for some $\vertexone\in \abag_j$.
\item \label{sequence-relabelinn-IntroEdge}
%$(\abag_i,\witnessset_i)=$ \\
$(\relabelingfunction_i(\abag_j), \action_{\dpcore}^{\treewidthvalue}(\relabelingfunction_i,\introedgegeneric{\vertexone}{\vertextwo}(\witnessset_j)))$
%\\
 for some distinct $\vertexone,\vertextwo\in \abag_j$.
	\item \label{sequence-relabeling-Join}
	%$(\abag_i,\witnessset_i) = $ \\
	$(\relabelingfunction_i(\abag_j), \action_{\dpcore}^{\treewidthvalue}(\relabelingfunction_i,\jointype(\witnessset_j,\action_{\dpcore}^{\treewidthvalue}(\pi,\witnessset_l))))$
	%\\
	for some $l\in \{0\}\cup[i-1]$ with $\abag_l=\abag_j$ and some $\pi\in \mathrm{Perm}(\abag_j)$.

		\end{enumerate}
		\end{enumerate}
		\end{definition}

		\begin{restatable}{lemma}{treedecexistssemiref}
		\label{lemma:tree decomposition-existence-sequence-relabel-semi-refutation}
		Let $\dpcore$ be a DP-core with identity cleaning and equipped with a witness action, and let $\relabelsequence=(\relabelingfunction_1,\dots,\relabelingfunction_m)$ be a sequence of $k$-relabeling functions. Let  $\dprefutation =(\abag_0,\witnessset_0)(\abag_1,\witnessset_1)\dots(\abag_m,\witnessset_m)$ be a semi $\relabelsequence$-relabeled $(\treewidthvalue,\dpcore)$-refutation, and $g:\abag_m\to [\treewidthvalue+1]$ be a $\treewidthvalue$-relabeling.
	Then, there is a $\treewidthvalue$-instructive tree decomposition $\abstractdecomposition_{\dprefutation}\in \allabstractdecompositionstreewidth{\treewidthvalue}$ and a function
	$\subtermToBag: \allsubterms(\abstractdecomposition_{\dprefutation}) \to \finitepowerset{\dpcore[\treewidthvalue].\allwitnesses}$
	such that the following conditions are satisfied for each subterm $\abstractdecomposition$ of $\abstractdecomposition_{\dprefutation}$.
\begin{enumerate}
    \item\label{existenceZero} $\topbag{\abstractdecomposition_{\dprefutation}} = g(\abag_m)$.
    \item\label{existenceOne} If $\abstractdecomposition=\abstractdecomposition_{\dprefutation}$, then $\subtermToBag(\abstractdecomposition) = \action_{\dpcore}^{\treewidthvalue}(g,\witnessset_m)$.
    \item\label{existenceTwo} if $\abstractdecomposition = \leaftype$, then
    $\subtermToBag(\abstractdecomposition) = \dpcore[\treewidthvalue].\initialsetgeneric$.
    \item\label{existenceThree} if $\abstractdecomposition = \introvertextype{\vertexone}(\abstractdecomposition_1)$ for some subterm $\abstractdecomposition_1$, then
    $\subtermToBag(\abstractdecomposition) = \introvertextype{\vertexone}(\subtermToBag(\abstractdecomposition_1))$.
    \item\label{existenceFour} if $\abstractdecomposition = \forgetvertextype{\vertexone}(\abstractdecomposition_1)$ for some subterm $\abstractdecomposition_1$, then
    $\subtermToBag(\abstractdecomposition) = \forgetvertextype{\vertexone}(\subtermToBag(\abstractdecomposition_1))$
    \item\label{existenceFive} if $\abstractdecomposition = \introedgegeneric{\vertexone}{\vertextwo}(\abstractdecomposition_1)$ for some subterm $\abstractdecomposition_1$, then
    $\subtermToBag(\abstractdecomposition) = \introedgegeneric{\vertexone}{\vertextwo}(\subtermToBag(\abstractdecomposition_1))$.
    \item\label{existenceSix} if $\abstractdecomposition = \jointype(\abstractdecomposition_1,\abstractdecomposition_2)$ for some subterms $\abstractdecomposition_1$ and $\abstractdecomposition_2$, then $\subtermToBag(\abstractdecomposition) = \jointype(\subtermToBag(\abstractdecomposition_1),\subtermToBag(\abstractdecomposition_2))$.
\end{enumerate}
\end{restatable}

	\begin{proof}
	We prove the lemma by induction on $m$.

Base case ($m=0$): Then $\dprefutation=(\abag_0,\witnessset_0)$ and by Condition~\ref{sequence-relabeling-condition-one} we have $(\abag_0,\witnessset_0)=(\emptyset,\dpcore[\treewidthvalue].\initialsetgeneric)$.
Let $\abstractdecomposition_{\dprefutation}\defeq \leaftype$ and set $\subtermToBag(\leaftype)\defeq \dpcore[\treewidthvalue].\initialsetgeneric$.
Since $\abag_0=\emptyset$, the relabeling $g$ is the unique empty function. For every witness $\awitness$ with $\mathrm{Lbl}_{\treewidthvalue}(\awitness)=\emptyset$, applying Definition~\ref{definition:witnessaction}.\ref{condition:relabeldpcore-condition-two} and Definition~\ref{definition:witnessaction}.\ref{condition:relabeldpcore-condition-three} yields $\action_{\dpcore}^{\treewidthvalue}(g,\awitness)=\awitness$; hence (by pointwise extension) $\action_{\dpcore}^{\treewidthvalue}(g,\witnessset_0)=\witnessset_0$.
Moreover, $\topbag{\leaftype}=\emptyset=g(\abag_0)$, so Condition~\ref{existenceZero} holds. Therefore Conditions~\ref{existenceZero}--\ref{existenceSix} hold.

Induction step: Assume $m\ge 1$ and that the statement holds for all semi-relabeled refutations of length $<m$.
Let $\dprefutation=(\abag_0,\witnessset_0)\dots(\abag_m,\witnessset_m)$ be a semi $\relabelsequence$-relabeled $(\treewidthvalue,\dpcore)$-refutation and fix $g:\abag_m\to [\treewidthvalue+1]$.
By Condition~\ref{sequence-relabeling-condition-three} of Definition~\ref{definition:RelabeledRefutation} (applied to $i=m$), there exists $j\in \{0\}\cup[m-1]$ such that $(\abag_m,\witnessset_m)$ is obtained from $(\abag_j,\witnessset_j)$ by one of the rules (and in the join case also some $l\in \{0\}\cup[m-1]$). We treat the cases.

\begin{enumerate}
\item \textbf{Introduce vertex.}
Assume
\[
(\abag_m,\witnessset_m)
=
\bigl(\relabelingfunction_m(\abag_j\cup\{\vertexone\}),\,\action_{\dpcore}^{\treewidthvalue}(\relabelingfunction_m,\introvertextype{\vertexone}(\witnessset_j))\bigr)
	\]
	for some $\vertexone\notin \abag_j$.
	Let
	\[
	h \defeq g\circ \bigl(\relabelingfunction_m|_{\abag_j\cup\{\vertexone\}}\bigr),
	\]
	which is well-defined because $\relabelingfunction_m(\abag_j\cup\{\vertexone\})=\abag_m=\domain(g)$.
	Apply the induction hypothesis to the prefix $(\abag_0,\witnessset_0)\dots(\abag_j,\witnessset_j)$ with relabeling $h|_{\abag_j}$.
	This yields $\abstractdecomposition_j$ and $\subtermToBag_j$ satisfying Conditions~\ref{existenceZero}--\ref{existenceSix} and such that
	\[
	\subtermToBag_j(\abstractdecomposition_j)=\action_{\dpcore}^{\treewidthvalue}(h|_{\abag_j},\witnessset_j).
	\]
	Define $\abstractdecomposition_{\dprefutation}\defeq \introvertextype{h(\vertexone)}(\abstractdecomposition_j)$.
	Extend $\subtermToBag_j$ by setting
	\[
	\subtermToBag(\abstractdecomposition_{\dprefutation})\defeq \introvertextype{h(\vertexone)}(\subtermToBag_j(\abstractdecomposition_j)).
	\]
	By the induction hypothesis, $\abstractdecomposition_j\in \allabstractdecompositionstreewidth{\treewidthvalue}$ and $\topbag{\abstractdecomposition_j}=h(\abag_j)$ (Condition~\ref{existenceZero} applied to the prefix).
	Since $\vertexone\notin \abag_j$ and $h$ is injective, we have
	\[
	h(\vertexone)\notin \topbag{\abstractdecomposition_j}.
	\]
	Hence $\introvertextype{h(\vertexone)}(\abstractdecomposition_j)\in \allabstractdecompositionstreewidth{\treewidthvalue}$ and
	\[
	\topbag{\abstractdecomposition_{\dprefutation}}
	= h(\abag_j\cup\{\vertexone\})
	= g(\abag_m).
	\]
		Thus Condition~\ref{existenceZero} holds, and Conditions~\ref{existenceTwo}--\ref{existenceSix} are immediate.
		For Condition~\ref{existenceOne}, since the refutation is well-formed (Condition~\ref{sequence-relabeling-condition-wellformed}), $\mathrm{Lbl}_{\treewidthvalue}(\witnessset_j)\subseteq \abag_j$.
		Therefore Definition~\ref{definition:witnessaction}.\ref{condition:relabeldpcore-condition-extension-invariance} gives
		\[
		\action_{\dpcore}^{\treewidthvalue}(h|_{\abag_j},\witnessset_j)=\action_{\dpcore}^{\treewidthvalue}(h,\witnessset_j).
		\]
		Let $S^\star\defeq \introvertextype{\vertexone}(\witnessset_j)$.
		Since $(\abag_m,\witnessset_m)$ is well-formed and $\abag_m=\relabelingfunction_m(\abag_j\cup\{\vertexone\})$, we have $\mathrm{Lbl}_{\treewidthvalue}(\witnessset_m)\subseteq \relabelingfunction_m(\abag_j\cup\{\vertexone\})$.
	Moreover, $\witnessset_m=\action_{\dpcore}^{\treewidthvalue}(\relabelingfunction_m,S^\star)$ by assumption, and this action application is well-defined by Condition~\ref{sequence-relabeling-condition-three}; hence label support transport (Definition~\ref{definition:witnessaction}.\ref{condition:relabeldpcore-condition-support-transport}, applied pointwise to sets) yields
	\[
	\mathrm{Lbl}_{\treewidthvalue}(\witnessset_m)=\relabelingfunction_m(\mathrm{Lbl}_{\treewidthvalue}(S^\star)).
	\]
	Therefore $\relabelingfunction_m(\mathrm{Lbl}_{\treewidthvalue}(S^\star))\subseteq \relabelingfunction_m(\abag_j\cup\{\vertexone\})$, and since $\relabelingfunction_m$ is injective we conclude $\mathrm{Lbl}_{\treewidthvalue}(S^\star)\subseteq \abag_j\cup\{\vertexone\}$ by Lemma~\ref{lemma:injective-image-containment}.
	Hence extension invariance yields $\action_{\dpcore}^{\treewidthvalue}(\relabelingfunction_m|_{\abag_j\cup\{\vertexone\}},S^\star)=\action_{\dpcore}^{\treewidthvalue}(\relabelingfunction_m,S^\star)=\witnessset_m$, and
		\[
		\begin{aligned}
		\subtermToBag(\abstractdecomposition_{\dprefutation})
		&= \introvertextype{h(\vertexone)}(\action_{\dpcore}^{\treewidthvalue}(h,\witnessset_j))\\
		&= \action_{\dpcore}^{\treewidthvalue}(h,S^\star)\\
		&= \action_{\dpcore}^{\treewidthvalue}\bigl(g,\,\action_{\dpcore}^{\treewidthvalue}(\relabelingfunction_m|_{\abag_j\cup\{\vertexone\}},S^\star)\bigr)\\
		&= \action_{\dpcore}^{\treewidthvalue}(g,\witnessset_m),
		\end{aligned}
		\]
	where the second equality follows from Definition~\ref{definition:witnessaction}.\ref{condition:relabeldpcore-condition-four} and the third from Definition~\ref{definition:witnessaction}.\ref{condition:relabeldpcore-condition-three}.

\item \textbf{Forget vertex.}
Assume $(\abag_m,\witnessset_m)=(\relabelingfunction_m(\abag_j\setminus\{\vertexone\}),\,\action_{\dpcore}^{\treewidthvalue}(\relabelingfunction_m,\forgetvertextype{\vertexone}(\witnessset_j)))$ for some $\vertexone\in \abag_j$.
Let $h\defeq g\circ \relabelingfunction_m|_{\abag_j\setminus\{\vertexone\}}$, which is well-defined since $\relabelingfunction_m(\abag_j\setminus\{\vertexone\})=\abag_m$.
Choose any injective extension $\hat h:\abag_j\to[\treewidthvalue+1]$ of $h$ (which exists by Lemma~\ref{lemma:injective-extension}).
Apply the induction hypothesis to the prefix ending at $j$ with relabeling $\hat h$ to obtain $\abstractdecomposition_j$ and $\subtermToBag_j$ with $\subtermToBag_j(\abstractdecomposition_j)=\action_{\dpcore}^{\treewidthvalue}(\hat h,\witnessset_j)$.
Define $\abstractdecomposition_{\dprefutation}\defeq \forgetvertextype{\hat h(\vertexone)}(\abstractdecomposition_j)$ and extend $\subtermToBag_j$ by setting $\subtermToBag(\abstractdecomposition_{\dprefutation})\defeq \forgetvertextype{\hat h(\vertexone)}(\subtermToBag_j(\abstractdecomposition_j))$.
By the induction hypothesis, $\abstractdecomposition_j\in \allabstractdecompositionstreewidth{\treewidthvalue}$ and $\topbag{\abstractdecomposition_j}=\hat h(\abag_j)$.
Since $\vertexone\in \abag_j$, we have $\hat h(\vertexone)\in \topbag{\abstractdecomposition_j}$, so $\forgetvertextype{\hat h(\vertexone)}(\abstractdecomposition_j)\in \allabstractdecompositionstreewidth{\treewidthvalue}$ and
\[
\topbag{\abstractdecomposition_{\dprefutation}}
= \hat h(\abag_j\setminus\{\vertexone\})
= h(\abag_j\setminus\{\vertexone\})
= g(\abag_m).
\]
Thus Condition~\ref{existenceZero} holds, and Conditions~\ref{existenceTwo}--\ref{existenceSix} are immediate.
For Condition~\ref{existenceOne}, Definition~\ref{definition:witnessaction}.\ref{condition:relabeldpcore-condition-five} implies
\[
	\subtermToBag(\abstractdecomposition_{\dprefutation})
	= \action_{\dpcore}^{\treewidthvalue}(\hat h,\forgetvertextype{\vertexone}(\witnessset_j)).
	\]
	Let $S^\star\defeq \forgetvertextype{\vertexone}(\witnessset_j)$.
	Since $(\abag_m,\witnessset_m)$ is well-formed (Condition~\ref{sequence-relabeling-condition-wellformed}) and $\abag_m=\relabelingfunction_m(\abag_j\setminus\{\vertexone\})$, we have
	\[
	\mathrm{Lbl}_{\treewidthvalue}(\witnessset_m)\subseteq \relabelingfunction_m(\abag_j\setminus\{\vertexone\}).
	\]
	Moreover, $\witnessset_m=\action_{\dpcore}^{\treewidthvalue}(\relabelingfunction_m,S^\star)$ by assumption, and this action application is well-defined by Condition~\ref{sequence-relabeling-condition-three}; hence label support transport (Definition~\ref{definition:witnessaction}.\ref{condition:relabeldpcore-condition-support-transport}, applied pointwise to sets) yields
	\[
	\mathrm{Lbl}_{\treewidthvalue}(\witnessset_m)=\relabelingfunction_m(\mathrm{Lbl}_{\treewidthvalue}(S^\star)).
	\]
	Therefore $\relabelingfunction_m(\mathrm{Lbl}_{\treewidthvalue}(S^\star))\subseteq \relabelingfunction_m(\abag_j\setminus\{\vertexone\})$, and since $\relabelingfunction_m$ is injective we conclude $\mathrm{Lbl}_{\treewidthvalue}(S^\star)\subseteq \abag_j\setminus\{\vertexone\}=\domain(h)$ by Lemma~\ref{lemma:injective-image-containment}.
	Hence extension invariance yields $\action_{\dpcore}^{\treewidthvalue}(\hat h,S^\star)=\action_{\dpcore}^{\treewidthvalue}(h,S^\star)$, and Definition~\ref{definition:witnessaction}.\ref{condition:relabeldpcore-condition-three} gives
		\[
		\begin{aligned}
		\subtermToBag(\abstractdecomposition_{\dprefutation})
		&= \action_{\dpcore}^{\treewidthvalue}(\hat h,S^\star)\\
		&= \action_{\dpcore}^{\treewidthvalue}(h,S^\star)\\
		&= \action_{\dpcore}^{\treewidthvalue}\bigl(g,\action_{\dpcore}^{\treewidthvalue}(\relabelingfunction_m,S^\star)\bigr)\\
		&= \action_{\dpcore}^{\treewidthvalue}(g,\witnessset_m).
		\end{aligned}
		\]

\item \textbf{Introduce edge.}
Assume
\[
(\abag_m,\witnessset_m)
=
\bigl(\relabelingfunction_m(\abag_j),\,\action_{\dpcore}^{\treewidthvalue}(\relabelingfunction_m,\introedgegeneric{\vertexone}{\vertextwo}(\witnessset_j))\bigr)
\]
for some distinct $\vertexone,\vertextwo\in \abag_j$.
Let $h\defeq g\circ \relabelingfunction_m|_{\abag_j}$ and apply the induction hypothesis to the prefix ending at $j$ with relabeling $h$, obtaining $\abstractdecomposition_j$ and $\subtermToBag_j$ with $\subtermToBag_j(\abstractdecomposition_j)=\action_{\dpcore}^{\treewidthvalue}(h,\witnessset_j)$.
Define $\abstractdecomposition_{\dprefutation}\defeq \introedgegeneric{h(\vertexone)}{h(\vertextwo)}(\abstractdecomposition_j)$ and set $\subtermToBag(\abstractdecomposition_{\dprefutation})\defeq \introedgegeneric{h(\vertexone)}{h(\vertextwo)}(\subtermToBag_j(\abstractdecomposition_j))$.
By the induction hypothesis, $\abstractdecomposition_j\in \allabstractdecompositionstreewidth{\treewidthvalue}$ and $\topbag{\abstractdecomposition_j}=h(\abag_j)$.
Since $\vertexone,\vertextwo\in \abag_j$, we have $h(\vertexone),h(\vertextwo)\in \topbag{\abstractdecomposition_j}$, so $\introedgegeneric{h(\vertexone)}{h(\vertextwo)}(\abstractdecomposition_j)\in \allabstractdecompositionstreewidth{\treewidthvalue}$ and $\topbag{\abstractdecomposition_{\dprefutation}}=\topbag{\abstractdecomposition_j}=h(\abag_j)=g(\abag_m)$.
Thus Condition~\ref{existenceZero} holds.
Let $S^\star\defeq \introedgegeneric{\vertexone}{\vertextwo}(\witnessset_j)$.
Since $(\abag_m,\witnessset_m)$ is well-formed and $\abag_m=\relabelingfunction_m(\abag_j)$, we have
\[
\mathrm{Lbl}_{\treewidthvalue}(\witnessset_m)\subseteq \relabelingfunction_m(\abag_j).
\]
Moreover, $\witnessset_m=\action_{\dpcore}^{\treewidthvalue}(\relabelingfunction_m,S^\star)$, so label support transport gives
\[
\mathrm{Lbl}_{\treewidthvalue}(\witnessset_m)=\relabelingfunction_m(\mathrm{Lbl}_{\treewidthvalue}(S^\star)).
\]
By Lemma~\ref{lemma:injective-image-containment}, $\mathrm{Lbl}_{\treewidthvalue}(S^\star)\subseteq \abag_j=\domain(h)$.
Hence
\[
\begin{aligned}
\subtermToBag(\abstractdecomposition_{\dprefutation})
&= \introedgegeneric{h(\vertexone)}{h(\vertextwo)}(\action_{\dpcore}^{\treewidthvalue}(h,\witnessset_j))\\
&= \action_{\dpcore}^{\treewidthvalue}(h,S^\star)\\
&= \action_{\dpcore}^{\treewidthvalue}\bigl(g,\action_{\dpcore}^{\treewidthvalue}(\relabelingfunction_m|_{\abag_j},S^\star)\bigr)\\
&= \action_{\dpcore}^{\treewidthvalue}(g,\witnessset_m),
\end{aligned}
\]
where the second equality is the introduce-edge equivariance axiom and the third is composition. The remaining conditions are immediate.

\item \textbf{Join.}
Assume $(\abag_m,\witnessset_m)=(\relabelingfunction_m(\abag_j),\,\action_{\dpcore}^{\treewidthvalue}(\relabelingfunction_m,\jointype(\witnessset_j,\action_{\dpcore}^{\treewidthvalue}(\pi,\witnessset_l))))$ for some $l\in \{0\}\cup[m-1]$ with $\abag_l=\abag_j$ and $\pi\in\mathrm{Perm}(\abag_j)$.
Let $h\defeq g\circ \relabelingfunction_m|_{\abag_j}$.
Apply the induction hypothesis to the prefixes ending at $j$ and $l$ with relabelings $h$ and $h\circ \pi$, obtaining decompositions $\abstractdecomposition_j,\abstractdecomposition_l$ and maps $\subtermToBag_j,\subtermToBag_l$ with
\[
\subtermToBag_j(\abstractdecomposition_j)=\action_{\dpcore}^{\treewidthvalue}(h,\witnessset_j),
\qquad
\subtermToBag_l(\abstractdecomposition_l)=\action_{\dpcore}^{\treewidthvalue}(h\circ \pi,\witnessset_l).
\]
Define $\abstractdecomposition_{\dprefutation}\defeq \jointype(\abstractdecomposition_j,\abstractdecomposition_l)$, using fresh copies of the two child terms if necessary; combine $\subtermToBag_j$ and $\subtermToBag_l$ on the respective subterms, and set
\[
\subtermToBag(\abstractdecomposition_{\dprefutation}) \defeq \jointype(\subtermToBag_j(\abstractdecomposition_j),\subtermToBag_l(\abstractdecomposition_l)).
\]
By the induction hypothesis, $\abstractdecomposition_j,\abstractdecomposition_l\in \allabstractdecompositionstreewidth{\treewidthvalue}$, $\topbag{\abstractdecomposition_j}=h(\abag_j)$, and $\topbag{\abstractdecomposition_l}=(h\circ \pi)(\abag_l)=(h\circ \pi)(\abag_j)=h(\abag_j)$.
Hence $\jointype(\abstractdecomposition_j,\abstractdecomposition_l)\in \allabstractdecompositionstreewidth{\treewidthvalue}$ and $\topbag{\abstractdecomposition_{\dprefutation}}=h(\abag_j)=g(\abag_m)$, so Condition~\ref{existenceZero} holds.
Because $(\abag_l,\witnessset_l)$ is well-formed and $\pi$ is a permutation of $\abag_j=\abag_l$, label support transport gives
\[
\mathrm{Lbl}_{\treewidthvalue}(\action_{\dpcore}^{\treewidthvalue}(\pi,\witnessset_l))\subseteq \abag_j.
\]
Let $S^\star\defeq \jointype(\witnessset_j,\action_{\dpcore}^{\treewidthvalue}(\pi,\witnessset_l))$.
Since $(\abag_m,\witnessset_m)$ is well-formed, $\abag_m=\relabelingfunction_m(\abag_j)$, and $\witnessset_m=\action_{\dpcore}^{\treewidthvalue}(\relabelingfunction_m,S^\star)$, label support transport and Lemma~\ref{lemma:injective-image-containment} imply
\[
\mathrm{Lbl}_{\treewidthvalue}(S^\star)\subseteq \abag_j=\domain(h).
\]
Therefore
\[
\begin{aligned}
\subtermToBag(\abstractdecomposition_{\dprefutation})
&= \jointype(\action_{\dpcore}^{\treewidthvalue}(h,\witnessset_j),\action_{\dpcore}^{\treewidthvalue}(h\circ \pi,\witnessset_l))\\
&= \jointype(\action_{\dpcore}^{\treewidthvalue}(h,\witnessset_j),\action_{\dpcore}^{\treewidthvalue}(h,\action_{\dpcore}^{\treewidthvalue}(\pi,\witnessset_l)))\\
&= \action_{\dpcore}^{\treewidthvalue}(h,S^\star)\\
&= \action_{\dpcore}^{\treewidthvalue}\bigl(g,\action_{\dpcore}^{\treewidthvalue}(\relabelingfunction_m|_{\abag_j},S^\star)\bigr)\\
&= \action_{\dpcore}^{\treewidthvalue}(g,\witnessset_m),
\end{aligned}
\]
where the second equality uses composition, the third uses join equivariance, and the fourth uses composition again. This proves Condition~\ref{existenceOne}, and the remaining conditions are immediate.
\end{enumerate}
\end{proof}

\end{document}
